\subsubsection{identify} \label{sec:identify}

System identification is implemented by the operation \operator{identify}.
In particular, the static method \texttt{identify} takes trajectory data as an input and identifies a suitable dynamic model that fits the data as accurately as possible.

The syntax for the operation \texttt{identify} is:

\begin{equation*}
\begin{split}
	& \texttt{sys} = \texttt{identify}(\texttt{traj}), \\
	& \texttt{sys} = \texttt{identify}(\texttt{traj},\texttt{options}), \\
	& \texttt{sys} = \texttt{identify}(\texttt{x},\texttt{t}),\\
	& \texttt{sys} = \texttt{identify}(\texttt{x},\texttt{t},\texttt{options}),\\
	& \texttt{sys} = \texttt{identify}(\texttt{x},\texttt{t},\texttt{u}),\\
	& \texttt{sys} = \texttt{identify}(\texttt{x},\texttt{t},\texttt{u},\texttt{options}),
\end{split}
\end{equation*}
with input arguments

\begin{center}
\renewcommand{\arraystretch}{1.3}
\begin{tabular}[t]{l p{13cm} }
	$\bullet$~\texttt{traj} & object of class \texttt{trajectory} (see \cref{sec:trajectory}) storing the trajectory data. \\
	$\bullet$~\texttt{x} & cell-array where each entry stores the states of one trajectory as a matrix $X \in \mathbb{R}^{n \times N}$, with $N$ being the number of time steps. \\
	$\bullet$~\texttt{t} & cell-array where each entry stores the time points of one trajectory as a vector $t \in \mathbb{R}^{N}$, with $N$ being the number of time steps. \\
	$\bullet$~\texttt{u} & cell-array where each entry stores the inputs of one trajectory as a matrix $U \in \mathbb{R}^{m \times N-1}$, with $N$ being the number of time steps. \\
	$\bullet$~\texttt{options} & struct containing algorithm settings \\
	& \begin{tabular}[t]{l p{10cm}}	
	 	--~\texttt{.alg} & string specifying the algorithm that is used for identification. The available algorithms are Dynamic Mode Decomposition \cite{Proctor2016} \texttt{`dmd'}, gradient-based optimization \texttt{`opt'}, SINDy \cite{Brunton2016} \texttt{`sparse'} (nonlinear systems only), and identification via nested functions \texttt{`nested'} (nonlinear systems only). For hybrid systems, the algorithm in \cite{Kochdumper2025} is used. The default value is \texttt{`dmd'}. \\
	 	--~\texttt{.dt} & sampling time $\Delta t$ for discrete-time models (classes \texttt{linearSysDT} and \texttt{nonlinearSysDT} only). The default value is the average time step size of the provided trajectory data. \\
	 	--~\texttt{.basis} & string specifying the template/basis functions used for identifying nonlinear dynamics. The availablte templates are all quadratic terms \texttt{`quad'}, all cubic terms \texttt{`cubic'}, trigonometric functions \texttt{`trig'}, and all together \texttt{`all'}. The default value is \texttt{`trig'}. \\
	 	--~\texttt{.phi} & custom user-defined basis/template functions for identifying nonlinear dynamics specified as a function handle, e.g. \texttt{options.phi = @(x,u) [x(1);x(2)*u(1)]} (all algorithms except \texttt{`nested'}). \\
	 	--~\texttt{.inpAff} & boolean specifying if the identified nonlinear model should be input-affine or not. The default value is \texttt{true}.\\
	 \end{tabular}
\end{tabular}
\end{center}

and output arguments

\begin{center}
\renewcommand{\arraystretch}{1.3}
\begin{tabular}[t]{l p{13cm} }
	$\bullet$~\texttt{sys} & dynamic system defined by any of the classes in \cref{sec:continuousDynamics,sec:hybridDynamics}, e.g., \texttt{linearSys}, \texttt{hybridAutomaton}, etc. \\
\end{tabular}
\end{center}


Let us demonstrate the operation \texttt{identify} by an example:

\begin{center}
\begin{minipage}[t]{0.58\textwidth}
	\footnotesize
	\input{./MATLABcode/example_identify}
\end{minipage}
\begin{minipage}[t]{0.4\textwidth}
	\vspace{10pt}

	\begin{verbatim}
	Command Window:	
		
	sys.A =

   -0.7034   -1.9941
    1.9745   -0.6953
	\end{verbatim}
\end{minipage}
\end{center}
